\begingroup

La présente annexe reprend, sous une forme consolidée, la chronologie du projet établie par Alain Fontan. Elle vise à documenter le cheminement ayant conduit le comité à formaliser le questionnaire comme instrument de cartographie des compétences, des disponibilités et des formes possibles de contribution des auditeurs. Les éléments ci-dessous conservent une fonction documentaire : ils éclairent la maturation de la démarche, sans se substituer à la démonstration principale du rapport.

\begin{figure}[H]
\centering
\resizebox{\linewidth}{!}{%
\begin{tikzpicture}[
  timeline/.style={-Latex, line width=0.8pt},
  event/.style={
    draw=black!65,
    fill=gray!8,
    rounded corners=1pt,
    align=left,
    font=\scriptsize,
    execute at begin node={\hyphenpenalty=10000\exhyphenpenalty=10000\emergencystretch=1em},
    text width=2.95cm,
    inner sep=3pt
  },
  dot/.style={circle, fill=black, inner sep=1.4pt}
]
\draw[timeline] (0,0) -- (15,0);
\foreach \x/\label in {
  0/{21/11/2025},
  4/{11/12/2025},
  8/{22/01/2026},
  12/{mars 2026}
}{
  \draw (\x,0.08) -- (\x,-0.08);
  \node[font=\scriptsize, below=0.18cm] at (\x,0) {\label};
}

\node[font=\bfseries\scriptsize] at (7.5,3.55) {Phase 1 -- cadrage et formalisation};

\node[dot] at (0.4,0) {};
\node[event, above=0.65cm] at (1.35,0) {Axes de travail : réserves, résilience locale, information, coordination et culture territoriale de sécurité globale.};

\node[dot] at (4.0,0) {};
\node[event, below=0.65cm] at (4.65,0) {Orientation vers le rôle des citoyens engagés et des associations régionales dans la prévention, l'accompagnement et la résilience.};

\node[dot] at (8.0,0) {};
\node[event, above=0.65cm] at (8.3,0) {Stabilisation de la problématique : contribution complémentaire, annuaire comme premier vivier, projet de questionnaire.};

\node[dot] at (10.8,0) {};
\node[event, below=0.65cm] at (10.8,0) {12/02--07/03/2026\\Confirmation d'une enquête rigoureuse ; présentation du cadre ; lettre de présentation à la présidence de l'AR 16.};

\node[dot] at (14.0,0) {};
\node[event, above=0.65cm] at (13.65,0) {19--20/03/2026\\Intérêt du bureau directeur ; souhait de prolongement vers 2027 ; échange méthodologique avec Victor Violier.};
\end{tikzpicture}
}

\vspace{1.1cm}

\resizebox{\linewidth}{!}{%
\begin{tikzpicture}[
  timeline/.style={-Latex, line width=0.8pt},
  event/.style={
    draw=black!65,
    fill=gray!8,
    rounded corners=1pt,
    align=left,
    font=\scriptsize,
    execute at begin node={\hyphenpenalty=10000\exhyphenpenalty=10000\emergencystretch=1em},
    text width=2.95cm,
    inner sep=3pt
  },
  dot/.style={circle, fill=black, inner sep=1.4pt}
]
\draw[timeline] (0,0) -- (15,0);
\foreach \x/\label in {
  0/{avr. 2026},
  4/{mai 2026},
  8/{juin 2026},
  14/{sept. 2027}
}{
  \draw (\x,0.08) -- (\x,-0.08);
  \node[font=\scriptsize, below=0.18cm] at (\x,0) {\label};
}

\node[font=\bfseries\scriptsize] at (7.5,3.55) {Phase 2 -- arbitrages et perspectives};

\node[dot] at (0.4,0) {};
\node[event, above=0.65cm] at (1.35,0) {03--22/04/2026\\Accord évoqué sur le projet, acculturation aux savoirs de crise, communication d'une ébauche du rapport.};

\node[dot] at (4.0,0) {};
\node[event, below=0.65cm] at (4.4,0) {05--10/05/2026\\Présentation d'une première version du rapport, puis information relative à l'invalidation du projet de questionnaire.};

\node[dot] at (8.0,0) {};
\node[event, above=0.65cm] at (8.35,0) {19/05, 21/05, 11/06/2026\\Analyse de l'annuaire 2022, échange sur la fin du projet d'enquête, derniers arbitrages.};

\node[dot] at (14.0,0) {};
\node[event, below=0.65cm] at (13.65,0) {Perspective de restitution ou de présentation institutionnelle, sous réserve de validation et de faisabilité.};
\end{tikzpicture}
}
\caption{Frise synthétique du cheminement du projet de questionnaire.}
\label{fig:frise_cheminement_questionnaire}
\end{figure}

\clearpage

\subsection*{Cheminement de l'élaboration du questionnaire}

\subsubsection*{Réunions et initiatives du comité}

\begingroup
\StyledTableSetup
\renewcommand{\arraystretch}{1.05}
\newcommand{\ChronoDate}[1]{\begin{minipage}[t]{\linewidth}\vspace{0pt}\RaggedRight #1\end{minipage}}
\newcommand{\ChronoItems}[1]{\begin{minipage}[t]{\linewidth}\vspace{0pt}\begin{itemize}[leftmargin=0.95em, nosep, topsep=0pt, partopsep=0pt, parsep=0pt, itemsep=0pt]#1\end{itemize}\end{minipage}}
\begin{xltabular}{\linewidth}{@{}>{\RaggedRight\arraybackslash}p{2.15cm}Y@{}}
\toprule
\textbf{Date} & \textbf{Éléments relevés dans la chronologie} \\
\midrule
\endfirsthead
\toprule
\textbf{Date} & \textbf{Éléments relevés dans la chronologie} \\
\midrule
\endhead
\midrule
\multicolumn{2}{r}{\small\itshape Suite page suivante} \\
\endfoot
\endlastfoot

\ChronoDate{21/11/2025} &
\ChronoItems{
    \item Identification d'axes de travail sur la place des réserves citoyennes et stratégiques.
    \item Rôle des compétences issues de l'IHEDN dans la résilience locale, la circulation de l'information et la sensibilisation de la population.
    \item Interactions entre associations régionales et autorités publiques.
    \item Besoin de coordination, de formation et de partage d'expérience.
    \item Renforcement d'une culture territoriale de sécurité globale.
} \\
\addlinespace

\ChronoDate{11/12/2025} &
\ChronoItems{
    \item Orientation de la réflexion vers le rôle que peuvent jouer les citoyens engagés et les associations régionales de l'IHEDN dans la prévention, la gestion et la résilience face aux crises majeures.
    \item Proposition d'une cartographie des compétences.
    \item Réflexion sur une éventuelle réserve IHEDN.
    \item Renforcement des actions de sensibilisation et de la coordination avec les autorités locales.
    \item Identification de points de vigilance et de défis de faisabilité et d'acceptabilité.
} \\
\addlinespace

\ChronoDate{22/01/2026} &
\ChronoItems{
    \item Formulation de la problématique du comité : les associations régionales de l'IHEDN n'ont pas vocation à décrire ou gérer directement une crise, mais à contribuer utilement à la préparation, à l'accompagnement des acteurs et à l'analyse des crises, en complément des autorités compétentes.
    \item Analyse de l'annuaire de l'\UnionIHEDN{} comme vivier significatif de compétences et d'expertises.
    \item Formalisation d'une proposition de questionnaire à destination des auditeurs de l'IHEDN.
    \item Objectif de cartographier les ressources disponibles, les domaines d'expertise et les profils volontaires pour intégrer un dispositif d'appui.
    \item Perspective d'un calendrier prévisionnel après consultation de l'\gls{irsem}.
    \item Intention de solliciter l'accord de la gouvernance de l'\ARSeize.
} \\
\addlinespace

\ChronoDate{12/02/2026} &
\ChronoItems{
    \item Confirmation de la proposition consistant à mieux connaître les membres de l'association par une enquête rigoureuse, crédible et structurée.
    \item Objectif de mieux connaître les membres de l'\ARSeize, de clarifier leur positionnement avant et après les crises, de proposer une intégration crédible dans les dispositifs existants et de renforcer la cohésion et la visibilité de l'association.
} \\
\addlinespace

\ChronoDate{03/03/2026} &
\ChronoItems{
    \item Rédaction et présentation de la problématique, du cadre général de l'enquête et de l'utilité du questionnaire.
} \\
\addlinespace

\ChronoDate{05/03/2026} &
\ChronoItems{
    \item Formulation d'une interrogation centrale sur la capacité réelle de mobilisation de l'\ARSeize{} en cas de crise majeure.
    \item Justification du questionnaire comme moyen de mieux connaître les profils, les compétences et la disponibilité des auditeurs.
    \item Mise en avant d'un possible effet de renforcement des liens entre les membres par la reconnaissance de leurs parcours et expertises.
    \item Inscription possible dans le contexte du cinquantième anniversaire de l'association et de l'\UnionIHEDN, ainsi que des 90 ans de l'IHEDN.
} \\
\addlinespace

\ChronoDate{06/03/2026} &
\ChronoItems{
    \item Envoi d'une lettre de présentation générale du projet à la présidence de l'\ARSeize\footnotemark.
} \\
\addlinespace
\midrule
\multicolumn{2}{r}{\small\itshape Suite page suivante} \\

\end{xltabular}
\endgroup
\footnotetext{Voir l'annexe \ref{annexe:lettre_CODIR}, qui reprend la lettre rédigée par Alain Fontan à destination du CODIR.}

\clearpage

\begingroup
\StyledTableSetup
\renewcommand{\arraystretch}{1.05}
\newcommand{\ChronoDate}[1]{\begin{minipage}[t]{\linewidth}\vspace{0pt}\RaggedRight #1\end{minipage}}
\newcommand{\ChronoItems}[1]{\begin{minipage}[t]{\linewidth}\vspace{0pt}\begin{itemize}[leftmargin=0.95em, nosep, topsep=0pt, partopsep=0pt, parsep=0pt, itemsep=0pt]#1\end{itemize}\end{minipage}}
\noindent\textit{Suite du tableau \ref{tab:reunions_initiatives_comite}.}

\begin{xltabular}{\linewidth}{@{}>{\RaggedRight\arraybackslash}p{2.15cm}Y@{}}
\toprule
\textbf{Date} & \textbf{Éléments relevés dans la chronologie} \\
\midrule
\endfirsthead
\toprule
\textbf{Date} & \textbf{Éléments relevés dans la chronologie} \\
\midrule
\endhead
\bottomrule
\endlastfoot

\ChronoDate{19/03/2026} &
\ChronoItems{
    \item Échange avec le bureau directeur de l'\ARSeize.
    \item Constat d'un intérêt pour la mise en place du questionnaire, celui-ci rejoignant une piste évoquée dans la feuille de route de l'\ARSeize.
    \item Convergence vers l'objectif de mieux structurer la contribution des auditeurs des associations en conciliant réalisme, disponibilité, compétences et cadre d'action.
    \item Expression d'un souhait de prolongement des travaux du comité jusqu'en 2027.
} \\
\addlinespace

\ChronoDate{09/04/2026} &
\ChronoItems{
    \item Échange sur la réalité de l'appartenance à la communauté des auditrices et auditeurs.
    \item Définition du périmètre concerné et insertion de cette réflexion dans la multiplicité des crises.
    \item Évocation d'un entretien informel avec la présidente de l'\UnionIHEDN, favorable à un questionnaire utile pour l'ensemble des associations, à tester à Paris--Île-de-France, dans une région rurale et en outre-mer.
} \\
\addlinespace

\ChronoDate{22/04/2026} &
\ChronoItems{
    \item Communication d'une ébauche du rapport du comité.
} \\
\addlinespace

\ChronoDate{05/05/2026} &
\ChronoItems{
    \item Présentation aux membres du comité de la première version du rapport intitulé « Crises majeures : quel rôle pour les citoyens engagés et les associations IHEDN ? ».
} \\
\addlinespace

\ChronoDate{19/05/2026} &
\ChronoItems{
    \item Partage d'une analyse de l'annuaire de l'\UnionIHEDN{} 2022 consacrée au vivier des auditeurs et aux limites du potentiel réellement mobilisable en cas de crise majeure.
    \item Approfondissement du rapport et clarification du positionnement de l'\ARSeize{} comme acteur de contribution indirecte mais stratégique, principalement en amont et en aval des crises, autour de la préparation, de la culture stratégique, de l'analyse et du retour d'expérience.
} \\
\addlinespace

\ChronoDate{21/05/2026} &
\ChronoItems{
    \item Échange autour de l'annonce de la fin du projet d'enquête.
} \\
\addlinespace

\ChronoDate{11/06/2026} &
\ChronoItems{
    \item Relecture et derniers arbitrages sur le contenu du rapport du comité avant son envoi au comité directeur de l'\ARSeize.
} \\

\end{xltabular}
\endgroup
\captionof{table}{Réunions et initiatives du comité dans le cheminement du projet de questionnaire.}
\label{tab:reunions_initiatives_comite}

\subsubsection*{Obstacles et conditions préalables identifiés}

\begin{itemize}[leftmargin=*]
    \item Cadre réglementaire : prise en compte du RGPD, applicable depuis le 25 mai 2018, et de la loi Informatique et Libertés du 6 janvier 1978 modifiée, mise en œuvre sous le contrôle de la CNIL.
    \item Préparation de l'envoi du questionnaire numérique, de sa présentation et de sa justification avant les étapes de relance et de collecte.
    \item Traitement et analyse des données, avec l'hypothèse d'une assistance institutionnelle ou universitaire.
    \item Prévision d'une restitution impliquant une organisation événementielle.
    \item Élaboration d'un calendrier prévisionnel allant de juin 2026 à septembre 2027\footnote{cf. figure \ref{fig:gantt-questionnaire-ar16}}.
\end{itemize}

\clearpage

\subsection*{Recherche de précédents et de ressources}

\subsubsection*{Études nationales et régionales consultées}

\begin{itemize}[leftmargin=*]
    \item Enquête nationale auprès des associations régionales de l'\UnionIHEDN{} de septembre 2015, consultée dans une double perspective : créer des synergies entre les composantes de l'\UnionIHEDN{} afin que l'efficacité de l'Union soit supérieure à la somme de celle de ses membres ; renforcer globalement la communauté des auditeurs sans méconnaître les spécificités des uns et des autres.
    \item Rapport du C2SD réalisé avec l'IEP de Toulouse et les AR IHEDN Toulouse et Montpellier, « L'esprit de défense au quotidien : les potentialités de développement de l'IHEDN », enquête anonyme auprès d'environ 300 acteurs locaux, par Jacques Aben, Marie-Dominique Charlier-Dagras et Jean-Pierre Marichy, 2000.
\end{itemize}

\subsubsection*{Ressources bibliographiques ou en ligne consultées}

\begin{itemize}[leftmargin=*]
    \item Archives de l'\UnionIHEDN, notamment la collection de la \textit{Revue Défense}.
    \item \textit{Sécurité, Défense, société civile : nouvelles solidarités}, Éditions Italiques, 2004.
    \item Vittorio Hösle, \textit{Les forces centrifuges planétaires. Une cartographie du temps présent à la lumière de la philosophie de l'histoire}, 2021, traduction Bruno Leray et Christine Griesmar.
    \item J. Sauvage, \textit{L'IHEDN : une vision globale de la politique de Défense de la France}, thèse d'histoire sous la direction de Maurice Waïsse, Reims, 1991.
    \item Robin Batard, \textit{Intégrer les contributions citoyennes aux dispositifs de gestion de crise : l'apport des médias sociaux}, thèse de l'Institut Polytechnique de Paris, 25 juin 2021.
    \item Rose Toki, \textit{Préparer les citoyens à la gestion des risques majeurs/risques de catastrophes à l'échelle des communes : une approche par l'instrumentation gestionnaire}, thèse de doctorat en sciences économiques et sciences de gestion, IAE Nantes, 13 décembre 2021.
    \item Colloque IHEDN au Sénat, \textit{Résilience d'une nation : y a-t-il une place pour le citoyen dans la sécurité de notre pays ?}, Lavauzelle, collection Les grands colloques de l'IHEDN, 2018.
    \item Semaine IHEDN Jeunes, AR Poitou-Charentes et Institut poitevin de recherche en sociologie de la connaissance, \textit{La personne citoyenne}, Poitiers, 31 août -- 5 septembre.
    \item \textit{Place Joffre : Cinquantenaire de l'IHEDN et des associations régionales d'auditeurs/IHEDN}, Éditions Hervas, 1999.
    \item Fondation pour les études de défense et Documentation française, \textit{La gestion de sortie de crise : actions civilo-militaires et opérations de reconstruction}, 1998.
    \item \textit{L'album des sessions : photos et liste des auditeurs des sessions nationales, régionales, européennes, africaines et malgaches}, 1948--1955, Éditions Hervas.
    \item Rapport du C2SD, IEP de Toulouse et AR IHEDN Toulouse et Montpellier, \textit{L'esprit de défense au quotidien : les potentialités de développement de l'IHEDN}, 2000.
    \item Les Jeunes IHEDN, GT OPS Jeux Olympiques et Paralympiques de Paris 2024, 2023.
    \item Sciences Po -- CrisisLab, Working Paper, « L'expérimentation générative et adaptative. Une approche de la préparation à la gestion de crise », 2025.
    \item Centre d'analyse comparative des systèmes politiques de l'Université Paris I, article « Les potentialités du développement des activités de l'IHEDN, enquête auprès des acteurs locaux », 1972.
    \item CNRS Info -- PEPR IRMa, entretien avec Soraya Boudia, « Il y a un besoin crucial d'innover en termes de politiques de gestion des risques et des crises », avril 2024.
    \item \gls{irsem}, Étude 70, Sarah Adjel, Angélique Palle et Noémie Rebière, \textit{Risques géopolitiques, crises et ressources naturelles -- Approches transversales et apport des sciences humaines}, septembre 2019.
    \item \gls{irsem}, Étude 95, \textit{L'armée, les Français et la crise sanitaire. Une enquête inédite}, juin 2022.
    \item CNRS, Lettre Sciences humaines et sociales, retour sur le séminaire « Les SHS à l'épreuve des risques et des catastrophes », 28 avril 2026.
    \item IHEMI, LIREC, Anne Muxel, Florian Opillard et Angélique Palle, « Comment développer la culture de crise ? », décembre 2023.
    \item IHEMI, LIREC, « Le volontariat et la mobilisation citoyenne : leviers de la gestion de crise ? », décembre 2025.
    \item Institut Montaigne, note d'enjeux, \textit{Esprit de défense : l'affaire de tous}, mars 2025.
    \item Lundi de l'IHEDN, « L'IHEDN au cœur d'un écosystème stratégique unique », mars 2026.
    \item \textit{Revue d'Anthropologie des connaissances}, article « Techno-politiques des vulnérabilités et production des catastrophes et crises », 20 janvier 2026.
    \item HAL open science, \textit{Bulletin of Sociological Methodology}, Marie Plessz, « Un protocole pour une enquête par questionnaire anonyme au sens du Règlement européen », 2020.
\end{itemize}

\subsubsection*{Rencontres d'acculturation à la notion de crise et aux méthodes d'enquête}

\begin{itemize}[leftmargin=*]
    \item Programme de recherche Risques/France 2030, « Les SHS à l'épreuve des risques et des catastrophes », Campus Saint-Germain, 6 novembre 2025.
    \item CNRS, projet RISC, « Risques et sociétés à l'ère des changements globaux : enjeux, savoirs et politiques », séance « Quels renouvellements des études de catastrophes ? Participation et mobilisation citoyenne dans la gestion des risques et des catastrophes », IPGP Jussieu, 17 décembre 2025.
    \item CNRS, Journées scientifiques annuelles PEPR Risques, « Construire des sciences des risques intégrées : défis et difficultés -- Recherche, expertise et action publique : quelles nouvelles perspectives pour une réduction durable des risques de catastrophes ? », 26--27 mars 2026.
    \item CNRS/EHESS, Campus Condorcet, « Savoirs et gouvernement des situations de crises et de catastrophes, XIXe--XXIe siècle : cadrage et instruments du gouvernement des catastrophes, mise en mémoire et oubli de l'histoire des catastrophes, politiques de l'alerte et de la réduction des catastrophes », 8 avril 2026.
    \item Université Paris I Panthéon-Sorbonne, journée d'étude « Protéger l'enquête : enjeux scientifiques et éthiques », groupe de recherche AFSP, avec notamment Victor Violier, Vanessa Frangville et Marielle Delbos, 20 mars 2026.
    \item Forum de Paris pour la défense et la stratégie, tables rondes « Engagement majeur, nouvelle génération », « Le rôle de la société civile dans la défense et la préservation de la paix » et « Préparer la nation au combat : les forces morales et la résilience de la nation », 24--26 mars 2025.
    \item Le Point -- Institut de France, Forum « Guerres et Paix », présentation de l'initiative « Au contact, citoyens, citoyennes » lancée par Muriel Domenach, ancienne ambassadrice de France auprès de l'OTAN, 1er avril 2026.
\end{itemize}

\clearpage

\subsection*{Cheminement de la faisabilité de la démarche}

\subsubsection*{Réassurance technique auprès d'universitaires}

\begingroup
\StyledTableSetup
\hyphenpenalty=10000
\exhyphenpenalty=10000
\emergencystretch=1em
\newcommand{\AlignedCell}[1]{\begin{minipage}[t]{\linewidth}\vspace{0pt}\RaggedRight #1\end{minipage}}
\newcommand{\AlignedBullets}[1]{\begin{minipage}[t]{\linewidth}\vspace{0pt}\begin{itemize}[leftmargin=1.15em, nosep, topsep=0pt, partopsep=0pt, parsep=0pt, itemsep=0pt]#1\end{itemize}\end{minipage}}
\begin{xltabular}{\linewidth}{@{}>{\RaggedRight\arraybackslash}p{3.6cm}YY@{}}
\toprule
\textbf{Interlocuteur} & \textbf{Points positifs} & \textbf{Points de vigilance} \\
\midrule
\endfirsthead
\toprule
\textbf{Interlocuteur} & \textbf{Points positifs} & \textbf{Points de vigilance} \\
\midrule
\endhead
\bottomrule
\endlastfoot

\AlignedCell{Victor Violier, sociologue à l'\gls{irsem}, pôle Défense et société\footnotemark} &
\AlignedBullets{
    \item Instrument de connaissance interne.
    \item Levier de structuration stratégique.
    \item Outil de valorisation institutionnelle.
    \item Fondement empirique pour toute proposition relative à l'intégration de l'\ARSeize{} dans une situation de crise majeure.
    \item Perspective d'une séance de restitution associant l'\ARSeize{} et l'\gls{irsem}.
} &
\AlignedBullets{
    \item Clarté de la problématique.
    \item Rigueur de la construction du questionnaire.
    \item Respect du cadre juridique applicable.
} \\
\addlinespace

\AlignedCell{Anne Muxel, directrice déléguée du \gls{cevipof}, directrice de recherches émérite en sociologie et sciences politiques au CNRS, ancienne directrice du domaine « Défense et société » de l'\gls{irsem}\footnotemark} &
\AlignedBullets{
    \item Démarche jugée probante.
    \item Intérêt de la structure et de la construction du questionnaire.
    \item Prise en compte de l'outre-mer dans le dispositif.
    \item Disponibilité pour accompagner la démarche.
} &
\AlignedBullets{
    \item RGPD et anonymisation du questionnaire à préserver.
} \\

\end{xltabular}
\endgroup
\footnotetext{Échanges avec des membres du comité le 20 février et le 20 mars 2026, dans le prolongement de la journée d'étude CNRS/EHESS/AFSP/Panthéon-Sorbonne « Protéger l'enquête : enjeux scientifiques et éthiques ».}
\captionof{table}{Échanges de réassurance technique auprès d'interlocuteurs universitaires.}
\label{tab:reassurance_technique_universitaires}

\clearpage

\subsubsection*{Réassurance interne par le biais d'échanges}

\begingroup
\StyledTableSetup
\hyphenpenalty=10000
\exhyphenpenalty=10000
\emergencystretch=1em
\newcommand{\AlignedCell}[1]{\begin{minipage}[t]{\linewidth}\vspace{0pt}\RaggedRight #1\end{minipage}}
\newcommand{\AlignedBullets}[1]{\begin{minipage}[t]{\linewidth}\vspace{0pt}\begin{itemize}[leftmargin=1.15em, nosep, topsep=0pt, partopsep=0pt, parsep=0pt, itemsep=0pt]#1\end{itemize}\end{minipage}}
\begin{xltabular}{\linewidth}{@{}>{\RaggedRight\arraybackslash}p{3.6cm}YY@{}}
\toprule
\textbf{Interlocuteur} & \textbf{Points positifs} & \textbf{Points de vigilance} \\
\midrule
\endfirsthead
\toprule
\textbf{Interlocuteur} & \textbf{Points positifs} & \textbf{Points de vigilance} \\
\midrule
\endhead
\bottomrule
\endlastfoot

\AlignedCell{Catherine Sarlandie de La Robertie, présidente de l'\UnionIHEDN, préfète, professeure des universités, rectrice d'académie, directrice de la formation continue Panthéon-Sorbonne} &
\AlignedBullets{
    \item Démarche pertinente à valoriser dans le cadre des anniversaires en cours.
    \item Dispositif à tester dans trois régions, dont une en outre-mer, avant un éventuel élargissement à l'ensemble des associations régionales.
    \item Proposition de présenter l'étude réalisée à la Sorbonne, Paris I, en septembre 2027.
} &
\AlignedBullets{
    \item Distinguer enquête et questionnaire.
    \item Discussion possible sur l'anonymisation des données.
} \\
\addlinespace

\AlignedCell{Pierre Vuillaume, général (2S), délégué général de l'\UnionIHEDN} &
\AlignedBullets{
    \item Intérêt pour la démarche.
    \item Caractère inédit de la proposition.
} &
\AlignedBullets{
    \item Ressources ou références historiques limitées pour documenter la préparation du projet.
    \item Dispositif considéré comme lourd.
} \\

\end{xltabular}
\endgroup
\footnotetext{\cite{muxel_cevipof}}
\captionof{table}{Échanges de réassurance interne sur la faisabilité et le positionnement du questionnaire.}
\label{tab:reassurance_interne_echanges}

Les conclusions de ces échanges ont été présentées en comité et ont fait l'objet de discussions.

\clearpage

\subsection*{Cheminement de la décision politique de l'\ARSeize}

\stepcounter{footnote}\xdef\noteLettreCODIRDecision{\number\value{footnote}}

\begin{figure}[H]
\centering
\resizebox{\linewidth}{!}{%
\begin{tikzpicture}[
  timeline/.style={-Latex, line width=0.8pt},
  event/.style={
    draw=black!65,
    fill=gray!8,
    rounded corners=1pt,
    align=left,
    font=\scriptsize,
    execute at begin node={\hyphenpenalty=10000\exhyphenpenalty=10000\emergencystretch=1em},
    text width=2.75cm,
    inner sep=3pt
  },
  dot/.style={circle, fill=black, inner sep=1.4pt}
]
\draw[timeline] (0,0) -- (15,0);
\foreach \x/\label in {
  0.7/{6 mars 2026},
  2.4/{7 mars 2026},
  6.0/{19 mars 2026},
  9.6/{3 avril 2026},
  14.0/{10 mai 2026}
}{
  \draw (\x,0.08) -- (\x,-0.08);
  \node[font=\scriptsize, below=0.18cm] at (\x,0) {\label};
  \node[dot] at (\x,0) {};
}

\node[event, above=0.65cm] at (0.9,0) {Lettre-courriel du comité d'études au bureau de l'\ARSeize{} proposant la mise en œuvre du questionnaire\footnotemark[\noteLettreCODIRDecision].};
\node[event, below=0.65cm] at (2.7,0) {Confirmation par la présidence de la transmission au bureau, avant discussion éventuelle en CODIR.};
\node[event, above=0.65cm] at (6.0,0) {Information donnée au comité sur l'intérêt suscité en CODIR ; aval officieux évoqué.};
\node[event, below=0.65cm] at (9.6,0) {Échange entre les présidentes : accord obtenu pour le projet de questionnaire.};
\node[event, above=0.65cm] at (14.0,0) {Information transmise à la présidente du comité : invalidation du projet de questionnaire.};
\end{tikzpicture}
}
\caption{Frise chronologique du cheminement de la décision politique relative au projet de questionnaire.}
\label{fig:decision_politique_questionnaire}
\end{figure}

\begingroup
\StyledTableSetup
\begin{xltabular}{\linewidth}{@{}p{2.6cm}Y@{}}
\toprule
\textbf{Date} & \textbf{Événement} \\
\midrule
\endfirsthead
\toprule
\textbf{Date} & \textbf{Événement} \\
\midrule
\endhead
\bottomrule
\endlastfoot

6 mars 2026 &
Le comité d'études adresse une lettre-courriel proposant le projet au bureau de l'\ARSeize{} pour une mise en œuvre éventuelle du questionnaire destiné à révéler une cartographie précise des ressources et compétences des membres de l'association\footnotemark[\noteLettreCODIRDecision]. \\
\addlinespace

7 mars 2026 &
La présidence de l'association régionale confirme la transmission du projet au bureau avant la poursuite éventuelle de la discussion lors du CODIR du 17 mars. \\
\addlinespace

19 mars 2026 &
Le comité d'études est informé de l'intérêt que suscite le projet de questionnaire au sein du CODIR ; un aval officieux du questionnaire est évoqué. \\
\addlinespace

3 avril 2026 &
Un échange entre la présidente de l'\ARSeize{} et la présidente du comité d'études permet d'obtenir un accord pour le projet de questionnaire. \\
\addlinespace

10 mai 2026 &
La présidente du comité est informée de l'invalidation du projet de questionnaire. \\

\end{xltabular}
\endgroup
\footnotetext[\noteLettreCODIRDecision]{Voir l'annexe \ref{annexe:lettre_CODIR}, qui reprend la lettre rédigée par Alain Fontan à destination du CODIR.}
\captionof{table}{Cheminement de la décision politique de l'\ARSeize{} relative au projet de questionnaire.}
\label{tab:decision_politique_questionnaire}

\endgroup
